\subsection*{只添加标签把请求历史缩成成员资格}

先限定一个简单用途：O1 的标签只表示累计发现的特征，添加以后一直保留。L 初始为空，A 观察到 gift，B 观察到 priority。双方各自保存集合，反复用并集合并，最后都是 $\{\mathrm{gift},\mathrm{priority}\}$。这里按不同标签值存储，十次重复设置 gift 只占一个成员。

相对于保存每次写入的完整请求和响应，这种结构确实少记了历史。它保留 O1、标签值和必要的来源验证，可以按订单与不同标签数计算空间。若业务还要知道是谁、何时、根据哪项授权发现 gift，就把这些依据纳入相应事实或审计记录，查询据此恢复成员来源。

产品现在要求移除 gift。假设 A 删除本地 gift，B 暂时离线，继续保存旧 $\{\mathrm{gift}\}$。恢复并集合并以后，gift 会被 B 带回来。集合中的缺失可能表示这里还没有收到添加，也可能表示这里已经删除，仅凭空集合无法区分这两段历史。

\CThirteenFigure{ch13-original08}{共同已有 gift，A 删除后再与 B 的旧副本取并集，标签重新出现}{fig:ch13-resurrection}

这与第十二章删除去重记录后的信息缺口相似。要区别已删除与尚未知，就给删除保留可传播依据。若产品规则允许删除后永远不能重加，可以分别保存已添加值与已删除值，用后者排除成员。O1 的礼物标签以后可能重新启用，本节采用能区分不同添加的方案。

\subsection*{删除源端已经看见的添加}

每次真正的标签添加取得唯一标记。例如 $(A,1)$ 表示来源 A 在这个对象上分配的第 1 次添加，简写为 A1。标记与标签值共同保存，gift 的一次添加就是 $(\mathrm{gift},A1)$。同一份添加效果重传沿用 A1，新添加使用新标记。

让 A、B 共同从已经包含 $(\mathrm{gift},A1)$ 的状态开始。A 移除 gift，记录删除集合 $\{A1\}$。B 与 A 断开以后，先移除自己看见的 A1，再重新添加 gift，取得 B1。合并时 A1 已删除，B1 保留，最终 gift 存在。

A 的移除只针对它已经观察到的添加。B1 在 A 做出移除时尚不可见，属于并发添加。这样的结构称为 observed-remove set，简称 OR-set，即已观察删除集合。对同一元素的添加与移除并发时，未被该移除观察到的添加保留，因此它也称添加优先集合~\cite{preguica2018crdt,bieniusa2012orset}。

\CThirteenFigure{ch13-original09}{A 移除已观察的 A1，B 重新添加取得 B1，合并后 B1 保留}{fig:ch13-observed-remove}

先用一份容易核对的表示保存两份只增长集合：$A_{\mathrm{add}}$ 保存已经见过的 $(\text{标签值},\text{唯一标记})$，$R_{\mathrm{remove}}$ 保存已删除标记。显示标签时，选取 $A_{\mathrm{add}}$ 中标记未进入 $R_{\mathrm{remove}}$ 的元素。添加扩大 $A_{\mathrm{add}}$，移除把源端当前可见的对应标记加入 $R_{\mathrm{remove}}$，合并对两份集合分别取并集。

这份内部状态按两份集合的包含关系增加，显示集合却可以变小。初态已有 A1，移除以后 $A_{\mathrm{add}}$ 保留 A1，$R_{\mathrm{remove}}$ 增加 A1，gift 从界面消失。随后旧添加 A1 又到，删除集合继续排除它。删除标记通常称为墓碑，作用是让旧副本携带的添加继续被识别为已删除。

若只保存 gift 已删除，又把所有未来 gift 都排除，B1 也会被压掉。若每次重试添加都新分配 B2、B3，OR-set 会把它们当成不同添加，后来移除时就可能要观察全部标记。业务若希望同一设置动作在整个恢复期只产生同一添加，就把该动作绑定到原标记，或者采用本来就表示单一持久事实的键。

OR-set 的并发语义是产品选择。它适合移除已经看见的那些添加，同时保留并发的新添加。若管理员撤销某项安全授权，要求并发授予也失效，就要采用覆盖这类并发的权限语义，并在对应授权仲裁中实施。

\subsection*{将逐项墓碑压成已观察范围}

前一表示在两份只增长集合中保留每次添加的标记。优化可以把已经见过且已删除的情况概括成因果上下文。上下文描述副本已经观察过哪些添加标记，配合当前还活着的标记，就能判断缺失是否表示删除。

先用显式集合表达这项信息。状态为 $(E,C)$：E 保存当前存活的 $(\text{标签值},\text{标记})$，C 保存已经观察的全部添加标记。添加一个新标记时，两者都加入相应内容。移除时只从 E 删掉已观察的活标记，C 保留它们。令 $\operatorname{dots}(E)$ 表示 E 中全部标记，则
\[
  R_{\mathrm{del}}=C-\operatorname{dots}(E)
\]
就是被上下文记住、当前已删除的那些标记。每个标记始终指向同一标签值，$\operatorname{dots}(E)\subseteq C$。

用 $(C,R_{\mathrm{del}})$ 可以核对这份信息的偏序：两份集合分别按包含比较。空初态中二者都为空。添加取得从未使用的新标记，扩大 C，$R_{\mathrm{del}}$ 保持。移除把已观察的活标记加入 $R_{\mathrm{del}}$，C 保持。两种迁移都膨胀，并保持 $R_{\mathrm{del}}\subseteq C$。

因此，直接用 E 的包含关系判断信息是否减少会出错。删除 gift 时，活项确实从 E 消失了，状态却增加了 A1 已经被删除的知识。将 $(E,C)$ 解码为 $(C,R_{\mathrm{del}})$，就能把这个变化放回同一信息偏序。

\CThirteenFigure{ch13-context-order}{E 可以因删除而缩小，已观察标记与已删除标记所表达的信息继续增加}{fig:ch13-context-order}

合并对 C、$R_{\mathrm{del}}$ 分别取并集，得到最小共同上界。再根据标记与标签值的固定关系，把 $C-R_{\mathrm{del}}$ 转回 E，就得到活项合并规则：双方都有的活标记保留。只在左边活着的标记，若右边尚未观察就保留，若右边已经观察却没有保留就按删除处理。右边独有的活标记作对称判断。

\begin{BookCode}
输入两份同对象状态 (E_left,C_left)、(E_right,C_right)
共同活项 = E_left 与 E_right 的交集
左侧新项 = E_left 中标记不属于 C_right 的项
右侧新项 = E_right 中标记不属于 C_left 的项
E = 对共同活项、左侧新项、右侧新项取并集
C = 对 C_left、C_right 取并集
返回 (E,C)
\end{BookCode}

下标 left、right 分别标记两份输入状态。第一行的共同活项取交集，保证双方共同保留的添加继续存在。两项新项再从单方活项中挑出对方尚未知晓的标记，已经在对方上下文中却没有活项的标记被排除。2012 年的优化报告采用这个判据，用每来源的连续计数概括已观察标记，并进一步合并同一元素、同一来源的冗余添加~\cite{bieniusa2012orset}。\footnote{2012 年报告将向量分量解释为已观察“n successive identifiers”，即从 1 开始连续的 n 个标记。它要求因果交付，脚注给出另记空洞的乱序扩展。本节的显式上下文便于展开空洞反例，原算法还按元素和来源合并活标记。}

\CThirteenFigure{ch13-original10}{活项 E 与上下文 C 共同区分尚未知和已删除，交集保留双方共同活项}{fig:ch13-context-merge}

一份上下文若包含 A1、A2、A3，可以用来源 A 的计数 3 表示这个连续前缀。多个来源的前缀合在一起，就是版本向量：每个分量说明对应来源的哪一段连续更新已经被观察。当前活项保留 A2，而向量说明已见 A1 至 A3，便可以推知 A1、A3 已经删除。逐项墓碑变成每源计数与活项的组合。

这个计数必须覆盖连续范围。设源 A 从空状态添加 gift，标记 A1，再添加 priority，标记 A2。B 只收到第二份添加效应器，第一份尚在路上。如果 B 直接令向量中 A 的分量为 2，却只保存活项 A2，那么它就声称已经观察 A1，并把缺失的 A1 解释成删除。A1 后来到达时，错误去重会丢掉它。

\CThirteenFigure{ch13-original11}{只收到 A2 时把连续前缀写成 2，会将尚未到达的 A1 误解释为已删除}{fig:ch13-vector-hole}

因果交付让 A1 先于 A2 到达，因而可以用连续向量。允许乱序时，可以保留已观察标记集合，或用连续前缀加例外集合记录空洞，补齐以后再压缩~\cite{bieniusa2012orset}。例如只见 A2 时保存前缀 0 和例外 $\{A2\}$，A1 到达后才压成前缀 2。第十二章的确认水位也只概括有依据的连续事实：版本向量记录已观察添加，确认水位记录已持久取得并确认的结果，各自根据对应事实推进。

墓碑或上下文的回收还要覆盖旧信息可能返回的路径。一种做法是在明确成员范围内，取得相关删除已被所有可继续同步成员持久观察的依据，再让落后快照、旧网络副本及退役来源经携带新上下文的重新加入流程恢复。若一台离线节点还能带着删除前状态直接参加合并，接收端就继续保存足以认出旧 A1 的删除依据。

本节显式 $(E,C)$ 先展开合并判据，C 按已观察标记增长，E 也可能为同一标签保留多次添加。原报告再把同一元素、同一来源的较早活标记合入较晚标记。例如 A 先后为 gift 产生 A1、A3，较晚添加承接同一来源的先前知识，优化结构只保留较晚的活标记 A3，向量保存已经见过的范围。来源 B 为 gift 产生的 B1 则有自己的分量与活标记。

\CThirteenFigure{ch13-orset-cost}{每个元素每个来源至多保留一个活标记，另用版本向量保存连续观察范围}{fig:ch13-orset-cost}

若当前有 n 个元素、固定 r 个来源，优化结构的活项至多为 $n\times r$，向量另有 r 个分量，按标记和计数分量的个数计为 $O(nr+r)$~\cite{bieniusa2012orset}。单个计数值增大时，其编码位数也会增加。来源动态增加时，向量分量随之增加，来源退出与重新加入规则也进入存储和恢复过程。

\subsection*{履约进度可以累计事实，再计算当前显示}

回到 O1 的查询投影。用履约事实集合 $\mathcal E_{\rm ful}$ 保存已核验的受理、发货和送达依据，每项都引用同一 O1/F1。受理事实引用履约库建立的真实 F1，即编号 501。发货事实引用原 \code{dispatch} 身份及下游确认，送达事实引用相应包裹与承运回执。对这条正常路径，可信的后阶段证据同时给出或允许核验前驱关系。

共同初态取 $\mathcal E_{\rm ful}$ 为空。投影先收到一份完整、已验证的发货证明，它包含 F1 的受理关系和原 \code{dispatch} 的成功依据，履约事实集合因而具有受理与发货事实，显示等级 2。较早的受理回报迟到，只会再次加入已有事实，等级保持 2。送达证据到来并验证以后，加入送达事实，显示等级 3。

若只有一个孤立的送达字样，投影先保存待核验来件，通过原 O1/F1、\code{dispatch} 和包裹关系补取依据。待核验输入与已确认 $\mathcal E_{\rm ful}$ 分开，显示从已确认集合计算，读到等级 3 就对应一条已经核实的业务事实链。

\CThirteenFigure{ch13-original12}{从已核验事实集合计算履约阶段，迟到回报补充已有事实，外部发货沿原身份执行}{fig:ch13-fulfillment}

这是本章额外设计的查询投影，增加字段与读取规则时要显式迁移。第十章 \code{orders} 表的 \code{fulfillment\_state} 当前只允许 \code{unfulfilled} 和 \code{fulfillment\_accepted} 两个值，直接写入 \code{dispatched} 会违反原检查约束。投影可以单独建对象，先启用证据验证与新读取器，再在需要时扩展原状态字段和约束。

同一订单在这条正常路径只有有限种阶段，重复通知可以汇到少量事实，接收端省去部分逐通知进度历史。订单数、F 编号、外部证据与包裹数继续决定总空间。每次通知的原结果回放、审计与问题追查若有独立要求，也保存对应记录。

取消、拦截、补发和退货进入以后，履约事实集合可以继续累计已提出取消、已发货、已批准补发等真实事实，显示规则按它们的具体关系判断。补发 AR1 形成新的履约身份，送达原包裹的证据只对应原包裹。集合合并让各副本逐步拥有相同事实，业务根据具体包裹、授权和前驱关系决定后续动作。
\FloatBarrier
